\chapter{Food Poisoning}

\section*{Definition}
Food poisoning is an illness caused by eating food or drinking water contaminated with bacteria, viruses, parasites, or their toxins. It typically causes nausea, vomiting, diarrhea, and abdominal cramps within hours to days after exposure.

\section*{What Happens in Your Body}
When you swallow contaminated food or water, the infectious organisms or their chemical toxins enter your digestive tract. Some organisms invade the lining of your stomach and intestines, causing inflammation and fluid leakage. Others release toxins that irritate the gut or are absorbed into your bloodstream, triggering vomiting and diarrhea as your body tries to expel the harmful substances. Most cases resolve on their own as your immune system clears the infection and your gut lining heals.

\section*{Causes}
\begin{itemize}
  \item Bacteria: Salmonella (poultry, eggs), E. coli (undercooked beef, contaminated produce), Campylobacter (raw or undercooked poultry), Staphylococcus aureus (improperly stored foods, creamy dishes), Listeria (deli meats, soft cheeses, unpasteurized dairy), Clostridium perfringens (reheated meat dishes), Vibrio (raw or undercooked shellfish)
  \item Viruses: Norovirus (contaminated food, water, or surfaces — very common), Hepatitis A, Rotavirus, Astrovirus
  \item Parasites: Giardia, Cryptosporidium, Toxoplasma (undercooked meat, contaminated water)
  \item Toxins: Pre-formed toxins from bacteria that cause rapid onset vomiting (Staph aureus, Bacillus cereus); toxic metals from acidic foods stored in certain containers
  \item Improper food handling: Undercooking, not refrigerating leftovers, cross-contamination between raw meat and ready-to-eat foods
\end{itemize}

\section*{When to See a Doctor}
See your doctor if:
\begin{itemize}
  \item Vomiting or diarrhea prevents you from keeping down any fluids for more than 24 hours
  \item You have signs of dehydration: Dry mouth, extreme thirst, little or no urination, dark urine, dizziness on standing
  \item Blood in your vomit or stool
  \item Fever above 38.5\textdegree{}C (101.3\textdegree{}F)
  \item Severe abdominal pain that does not come and go (cramping)
  \item Diarrhea lasts more than 3 days in adults
  \item You are pregnant, over 65, or have a weakened immune system
  \item The person who is ill is an infant or young child
  \item Vision changes, muscle weakness, or trouble speaking (rare — possible botulism, seek emergency care)
\end{itemize}

\section*{Related Symptoms}
\begin{itemize}
  \item Nausea
  \item Vomiting
  \item Diarrhea (watery or bloody)
  \item Abdominal cramps
  \item Fever and chills
  \item Headache
  \item Muscle aches
  \item Fatigue
  \item Loss of appetite
  \item Dehydration
\end{itemize}
\textbf{Important:} Food poisoning from different sources can look very similar. The specific cause is not always identified, but knowing what you ate and when others who shared your meal got sick can help narrow it down.

\section*{Self-Care / Home Management}
\begin{itemize}
  \item Rest your stomach — avoid solid food for a few hours and sip small amounts of clear liquids (water, clear broth, oral rehydration solution).
  \item Rehydrate gradually — take frequent small sips rather than large amounts at once.
  \item Oral rehydration solutions (Pedialyte, DripDrop, or a homemade mix of water, salt, and sugar) replace lost fluids and electrolytes better than water alone.
  \item When you feel ready, eat bland, easy-to-digest foods: Crackers, toast, rice, bananas, applesauce.
  \item Avoid dairy, fatty or spicy foods, caffeine, and alcohol until you are fully recovered.
  \item Do not take anti-diarrhea medications (loperamide / Imodium) if you have bloody diarrhea or a fever — they can keep harmful organisms in your body.
  \item Wash your hands frequently and avoid preparing food for others while you are sick.
\end{itemize}

\section*{How Your Doctor Will Evaluate You}
\begin{itemize}
  \item Your doctor will ask what you ate, where you ate it, when symptoms started, and whether anyone else who shared the meal is sick.
  \item They will check your vital signs and look for signs of dehydration.
  \item A stool sample may be tested for bacteria, viruses, or parasites if the illness is severe, prolonged, or part of a suspected outbreak.
  \item Blood tests may be ordered if you are severely dehydrated or if listeria or botulism is suspected.
  \item Imaging (CT scan of the abdomen) is sometimes needed if the pain is severe enough to suggest appendicitis or other surgical conditions.
\end{itemize}

\section*{Treatment Options}
\begin{itemize}
  \item Most cases resolve on their own within 1 to 3 days with rest and hydration.
  \item Antibiotics are prescribed only for specific bacterial infections (salmonella, listeria, campylobacter) that are severe or affect high-risk patients.
  \item Antiparasitic medications are used for parasitic infections like giardia.
  \item Hospitalization with intravenous fluids may be needed for severe dehydration.
  \item Probiotics may help restore healthy gut bacteria after a diarrheal illness, but evidence is limited.
  \item Prevention is key: Cook meat thoroughly, wash produce, refrigerate leftovers within 2 hours, and wash your hands before handling food.
\end{itemize}

\section*{References}
\begin{thebibliography}{99}
\bibitem{food_poisoning:ref1} Scallan E, Hoekstra RM, Angulo FJ, et al. ``Foodborne Illness Acquired in the United States: Major Pathogens.'' Emerg Infect Dis. 2011;17(1):7--15.
\bibitem{food_poisoning:ref2} DuPont HL. ``Clinical Practice: Bacterial Diarrhea.'' N Engl J Med. 2009;361(16):1560--1569.
\bibitem{food_poisoning:ref3} Tarr PI, Gordon CA, Chandler WL. ``Shiga-Toxin-Producing Escherichia coli and Haemolytic Uraemic Syndrome.'' Lancet. 2005;365(9464):1073--1086.
\bibitem{food_poisoning:ref4} Thielman NM, Guerrant RL. ``Clinical Practice: Acute Infectious Diarrhea.'' N Engl J Med. 2004;350(1):38--47.
\bibitem{food_poisoning:ref5} Shane AL, Mody RK, Crump JA, et al. ``2017 Infectious Diseases Society of America Clinical Practice Guidelines for the Diagnosis and Management of Infectious Diarrhea.'' Clin Infect Dis. 2017;65(12):e45--e80.
\bibitem{food_poisoning:ref6} Baird-Parker AC. ``Foodborne Illness: Foodborne Salmonellosis.'' Lancet. 1990;336(8725):1231--1235.
\end{thebibliography}
